\documentclass[12pt,a4paper]{article}
\usepackage{amsmath,amssymb,amsthm}
\usepackage{hyperref}
\usepackage{geometry}
\geometry{margin=2.5cm}
\usepackage{enumitem}
\usepackage{booktabs}

\newtheorem{theorem}{Theorem}[section]
\newtheorem{lemma}[theorem]{Lemma}
\newtheorem{corollary}[theorem]{Corollary}
\newtheorem{proposition}[theorem]{Proposition}
\theoremstyle{definition}
\newtheorem{definition}[theorem]{Definition}
\theoremstyle{remark}
\newtheorem{remark}[theorem]{Remark}

\title{Dirac vs.\ Majorana Neutrinos in Recognition Science:\\
Double-Entry Structure Determines the Dirac Nature}

\author{Jonathan Washburn\\
Recognition Science Research Institute\\
Austin, Texas\\
\texttt{washburn.jonathan@gmail.com}}

\date{March 2026}

\begin{document}
\maketitle

\begin{abstract}
We address the Dirac--Majorana question for neutrinos within
Recognition Science (RS).  The ledger's double-entry bookkeeping
constrains the neutrino mass term.  A Dirac mass
$\bar{\nu}_L m_D \nu_R$ requires a distinct right-handed partner
(credit entry for each debit entry).  A Majorana mass
$\bar{\nu}_L^c m_M \nu_L$ identifies the neutrino with its own
antiparticle (self-debit), violating double-entry bookkeeping.
RS determines that neutrinos are Dirac: every fermion entry on the
ledger requires a distinct antiparticle counterpart.  The
generation count $N_\nu = \mathrm{face\_pairs}(3) = 3$ is fixed
by the $D = 3$ cube geometry; a fourth generation is excluded.
The observational consequence: neutrinoless double-beta decay
($0\nu\beta\beta$) should not be observed.  We analyze the
implications for the Majorana phases in the PMNS matrix and
derive the lepton number conservation that follows.
\end{abstract}

%% ============================================================
\section{Recognition Science Framework}
\label{sec:framework}
%% ============================================================

The Dirac--Majorana question is resolved by the RS ledger axioms,
which uniquely determine the cost functional and the double-entry
structure of all particle interactions.

\begin{definition}[RS Cost Functional]
For $x > 0$: $J(x) = \tfrac{1}{2}(x + x^{-1}) - 1$.
\end{definition}

\noindent\textbf{Axiom A1 (Normalization):} $J(1) = 0$.\\
\textbf{Axiom A2 (Recognition Composition Law, RCL):}
$J(xy)+J(x/y)=2J(x)J(y)+2J(x)+2J(y)$ for all $x,y>0$.\\
\textbf{Axiom A3 (Calibration):} $J''_{\log}(0) = 1$, where
$J_{\log}(t) := J(e^t)$.

\begin{theorem}[Cost Uniqueness]
\label{thm:unique}
The unique continuous solution to \textup{(A1)--(A3)} is
$J(x) = \tfrac{1}{2}(x + x^{-1}) - 1$.
\end{theorem}

\begin{proof}[Proof sketch]
Set $H(t) = J_{\log}(t)+1$.  Axiom A2 transforms into the d'Alembert
equation $H(s+t)+H(s-t)=2H(s)H(t)$.  By the Acz\'el classification
theorem~\cite{Aczel1966}, the unique continuous solution with $H(0)=1$
and $H''(0)=1$ is $H(t)=\cosh t$, giving
$J(x)=\tfrac{1}{2}(x+x^{-1})-1$.
\end{proof}

\noindent The RS ledger implements a double-entry bookkeeping structure:
every debit (particle entry) requires a distinct credit (antiparticle
entry).  This structural postulate, together with $D=3$ (forced by
$2^D=8$), determines both the Dirac nature of neutrinos and the
generation count $N_\nu = \mathrm{face\_pairs}(3) = 3$.

%% ============================================================
\section{Introduction}
\label{sec:intro}
%% ============================================================

Whether neutrinos are Dirac particles (with distinct
antiparticles, $\nu \neq \bar{\nu}$) or Majorana particles
(their own antiparticles, $\nu = \bar{\nu}$) is one of the most
fundamental open questions in particle
physics~\cite{Bilenky1987}.  Recognition Science (RS) resolves
this via the double-entry bookkeeping axiom of the recognition
ledger~\cite{RS_Axioms}: every debit must have a distinct
credit entry, which forces a Dirac structure.  The generation
count is fixed by the $D = 3$ cube geometry.  The consequences
are sharp and falsifiable.
\begin{itemize}
  \item \textbf{Dirac}: Lepton number is conserved.
  Neutrinoless double-beta decay is forbidden.  The PMNS matrix
  has one Dirac CP phase.
  \item \textbf{Majorana}: Lepton number is violated by 2 units.
  $0\nu\beta\beta$ is allowed.  The PMNS matrix has two additional
  Majorana phases.
\end{itemize}

The Standard Model is agnostic.  RS is not.

%% ============================================================
\section{Double-Entry Constraint}
\label{sec:double-entry}
%% ============================================================

\begin{definition}[Double-Entry Bookkeeping]
The recognition ledger records every event as a debit-credit
pair: each particle entry (debit) has a distinct antiparticle
entry (credit).  The total charge (summed over all entries) is
zero.
\end{definition}

\begin{proposition}[RS Structural Prediction: Dirac Neutrinos]
\label{prop:dirac}
The double-entry bookkeeping postulate of the RS ledger requires
neutrinos to be Dirac particles.
\end{proposition}

\begin{proof}
The double-entry postulate states: every particle entry (debit) on
the recognition ledger must have a distinct antiparticle entry
(credit).  Consider the two possible neutrino mass terms:
\begin{enumerate}[label=(\roman*)]
  \item \textbf{Dirac}: $\mathcal{L}_D = -m_D\,
  \bar{\nu}_L \nu_R + \mathrm{h.c.}$  The left-handed neutrino
  $\nu_L$ is the debit entry; the right-handed neutrino $\nu_R$
  is the distinct credit entry.  This satisfies double-entry.

  \item \textbf{Majorana}: $\mathcal{L}_M = -\tfrac{1}{2}\,
  m_M\,\bar{\nu}_L^c \nu_L + \mathrm{h.c.}$  Here $\nu_L^c$ is
  identified with $\nu_L$ itself ($\nu = \bar{\nu}$).  The same
  ledger entry serves as both debit and credit---a ``self-debit''
  that violates the double-entry postulate.
\end{enumerate}
Since double-entry is a structural postulate of the RS ledger
(not derived from A1--A3 but from the bookkeeping axiom of the
recognition process), option~(ii) is excluded within RS.
\end{proof}

\begin{remark}[Status of the double-entry postulate]
The double-entry bookkeeping structure is a postulate of the RS
ledger architecture, logically independent of but consistent with
the cost axioms A1--A3.  The Dirac prediction is therefore
conditional on accepting this ledger postulate.  The prediction
is, however, sharply testable via $0\nu\beta\beta$ experiments.
\end{remark}

%% ============================================================
\section{Generation Count}
\label{sec:generations}
%% ============================================================

\begin{definition}[Face Pairs in $D = 3$]
A $D$-dimensional hypercube has $2D$ faces.  Its face-pair count
is $\mathrm{face\_pairs}(D) = D$, since each pair of parallel
faces constitutes one spatial direction.  For $D = 3$:
$\mathrm{face\_pairs}(3) = 3$.
\end{definition}

\begin{theorem}[Three Neutrino Generations]
\label{thm:three}
$N_\nu = \mathrm{face\_pairs}(3) = 3$: exactly three neutrino
generations ($\nu_e, \nu_\mu, \nu_\tau$).
\end{theorem}

\begin{theorem}[Fourth Generation Excluded]
\label{thm:no-fourth}
$\mathrm{face\_pairs}(3) = 3 \neq 4$: a fourth neutrino
generation does not exist.
\end{theorem}

\begin{remark}
The LEP measurement of the invisible $Z$ width gives
$N_\nu = 2.9840 \pm 0.0082$~\cite{LEP2006}, consistent with
exactly 3.
\end{remark}

%% ============================================================
\section{Lepton Number Conservation}
\label{sec:lepton}
%% ============================================================

\begin{corollary}[Lepton Number Conserved]
Since neutrinos are Dirac, lepton number $L$ is a conserved
quantum number.  Processes with $|\Delta L| = 2$ (such as
$0\nu\beta\beta$) are forbidden.
\end{corollary}

%% ============================================================
\section{PMNS Matrix Structure}
\label{sec:pmns}
%% ============================================================

\begin{theorem}[PMNS Phases]
For Dirac neutrinos, the PMNS mixing matrix has:
\begin{itemize}
  \item 3 mixing angles: $\theta_{12}, \theta_{23}, \theta_{13}$.
  \item 1 Dirac CP phase: $\delta_{CP}$.
  \item 0 Majorana phases (these are unphysical for Dirac
  neutrinos).
\end{itemize}
\end{theorem}

\begin{remark}
The RS prediction for the Dirac CP phase is
$\delta_{CP} = -\pi/2$ ($270°$), from the 8-tick
quarter-period (see companion paper).  Current
global fits favor $\delta_{CP} \approx 195°$--$255°$~\cite{PDG2022},
consistent at $\sim 2\sigma$.
\end{remark}

%% ============================================================
\section{Observational Tests}
\label{sec:tests}
%% ============================================================

\begin{enumerate}
  \item \textbf{$0\nu\beta\beta$ not observed}: If neutrinos
  are Dirac, neutrinoless double-beta decay is forbidden.
  Current best limits:
  $T_{1/2}^{0\nu}(^{136}\mathrm{Xe}) > 2.3 \times
  10^{26}\;\mathrm{yr}$ at 90\% CL
  (KamLAND-Zen 800)~\cite{KamLANDZen2022};
  $T_{1/2}^{0\nu}(^{76}\mathrm{Ge}) > 1.8 \times 10^{26}
  \;\mathrm{yr}$ (GERDA)~\cite{GERDA2020}.  RS predicts both
  will continue to grow without finding a signal.

  \item \textbf{Falsifier}: Observation of $0\nu\beta\beta$ at
  any rate consistent with a Majorana mass would falsify the RS
  double-entry prediction.

  \item \textbf{$N_\nu = 3$ exactly}: No sterile neutrinos
  mixing with active neutrinos.

  \item \textbf{Lepton number conserved}: No lepton-number
  violating processes at any energy.
\end{enumerate}

%% ============================================================
\section{Comparison}
\label{sec:comparison}
%% ============================================================

\begin{center}
\renewcommand{\arraystretch}{1.3}
\begin{tabular}{@{}lcc@{}}
\toprule
\textbf{Framework} & \textbf{Nature} &
\textbf{$0\nu\beta\beta$?} \\
\midrule
Standard Model (minimal) & Either & Unknown \\
See-saw (Type I) & Majorana & Yes \\
Left-right symmetric & Majorana & Yes \\
RS (double-entry) & Dirac & No \\
\bottomrule
\end{tabular}
\end{center}

%% ============================================================
\section{Conclusion}
\label{sec:conclusion}
%% ============================================================

Recognition Science's double-entry ledger structure determines
that neutrinos are Dirac particles: each neutrino has a distinct
antineutrino, and lepton number is conserved.  Three generations
follow from $D = 3$.  Neutrinoless double-beta decay is predicted
to be absent---a sharp, falsifiable prediction testable by
current and near-future experiments.

%% ============================================================
\begin{thebibliography}{99}

\bibitem{RS_Axioms}
J.~Washburn, ``The Algebra of Reality: A Recognition Science Derivation
of Physical Law,'' \emph{Axioms} \textbf{15}(2), 90 (2026).
\texttt{doi:10.3390/axioms15020090}

\bibitem{Aczel1966}
J.~Acz\'el, \emph{Lectures on Functional Equations and Their
Applications} (Academic Press, 1966).

\bibitem{Bilenky1987}
S.~M.~Bilenky and S.~T.~Petcov, ``Massive Neutrinos and
Neutrino Oscillations,'' Rev.\ Mod.\ Phys.\ \textbf{59}, 671
(1987).

\bibitem{LEP2006}
ALEPH, DELPHI, L3, OPAL, SLD Collaborations, ``Precision
Electroweak Measurements on the Z Resonance,'' Phys.\ Rep.\
\textbf{427}, 257 (2006).

\bibitem{PDG2022}
S.~Navas \textit{et al.}\ (Particle Data Group), ``Review of Particle
Physics,'' Phys.\ Rev.\ D \textbf{110}, 030001 (2024).

\bibitem{KamLANDZen2022}
KamLAND-Zen Collaboration (S.~Abe \textit{et al.}),
``First Search for the Majorana Nature of Neutrinos in the
Inverted Mass Ordering Region with KamLAND-Zen,''
Phys.\ Rev.\ Lett.\ \textbf{130}, 051801 (2023).

\bibitem{GERDA2020}
GERDA Collaboration, ``Final results of GERDA on the search for
neutrinoless double-$\beta$ decay,'' Phys.\ Rev.\ Lett.\
\textbf{125}, 252502 (2020).

\end{thebibliography}

\end{document}
